% ===========================================================================
% BACKUP SLIDES
% Statements copied from the report (FINAL/report), with the report numbering.
% Every frame has a Back button to the main slide that links to it.
% ===========================================================================

% ---------------------------------------------------------------------------
% Preliminaries
% ---------------------------------------------------------------------------

\begin{frame}[label=bk-prelim]{Matrices (Section 1.1.2)}
  \scriptsize
  \begin{block}{Definition 1.5 (Non-negative and stochastic matrices) [3]}
    A matrix $A = [a_{ij}] \in \mathbb{R}^{n \times n}$ is
    (i) \emph{non-negative} (\emph{positive}) if $a_{ij} \geq 0$ ($a_{ij} > 0$, resp.) for
    all $i, j \in \{1, \dots, n\}$;
    (ii) \emph{row-stochastic} if it is non-negative and $A \mathbf{1}_n = \mathbf{1}_n$;
    (iii) \emph{column-stochastic} if it is non-negative and
    $A^\top \mathbf{1}_n = \mathbf{1}_n$;
    (iv) \emph{doubly-stochastic} if it is row- and column-stochastic.
  \end{block}
  \begin{block}{Theorem 1.6 (Spectral properties of row-stochastic matrices) [3]}
    For a row-stochastic matrix $A \in \mathbb{R}^{n \times n}$
    (i) $1 \in \operatorname{spec}(A)$;
    (ii) $\operatorname{spec}(A)$ is a subset of the unit disk and $\rho(A) = 1$.
  \end{block}
  \begin{block}{Definition 1.7 (Semi-convergent and convergent matrices) [3]}
    A matrix $A \in \mathbb{R}^{n \times n}$ is
    (i) \emph{semi-convergent} if $\lim_{k \to +\infty} A^k$ exists (finite);
    (ii) \emph{convergent} (or \emph{Schur stable}) if it is semi-convergent and
    $\lim_{k \to +\infty} A^k = 0_{n \times n}$.
  \end{block}
  \detailbuttons{\go{bk-prelim-b}{Next}\ \ret{m-fdg-model}}
\end{frame}

\begin{frame}[label=bk-prelim-b]{Spectral radius and regularity}
  \scriptsize
  \begin{block}{Theorem 1.8 (Convergence and spectral radius) [3]}
    Consider a matrix $A \in \mathbb{R}^{n \times n}$ and let
    $\rho(A) := \max\{|\lambda| \mid \lambda \in \operatorname{spec}(A)\}$ (spectral radius
    of $A$). Then:
    (i) $A$ convergent $\iff$ $\rho(A) < 1$;
    (ii) $A$ semi-convergent and not convergent $\iff$ $1$ is a semisimple eigenvalue of
    $A$ and all the other eigenvalues have modulus strictly less than $1$.
  \end{block}
  \begin{block}{Definition 1.9 (Regularity) [5]}
    We call the matrix $A$ \emph{regular} if the limit $A^\infty = \lim_{k \to \infty} A^k$
    exists and \emph{fully regular} if, additionally, the rows of $A^\infty$ are identical
    (that is, $A^\infty = \mathbf{1}_n w^\top$ for some $w \in \mathbb{R}^n$).
  \end{block}
  \detailbuttons{\go{bk-prelim2}{Next}\ \ret{m-fdg-model}}
\end{frame}

\begin{frame}[label=bk-prelim2]{Primitive, substochastic}
  \scriptsize
  \begin{block}{Definition 1.10 (Irreducible and primitive matrices) [3]}
    A non-negative matrix $A \in \mathbb{R}^{n \times n}$, $n \geq 2$, is
    (i) \emph{irreducible} if $\sum_{k=0}^{n-1} A^k$ is a positive matrix;
    (ii) \emph{primitive} if there exists $k \in \mathbb{N}$ such that $A^k$ is positive;
    (iii) \emph{reducible} if it is not irreducible.
  \end{block}
  \begin{block}{Theorem 1.11 (Primitive matrices and digraphs) [2]}
    Let $G = (V, E, \{a_{ij}\}_{(i,j) \in E})$ be a weighted digraph of order $n \geq 2$
    with adjacency matrix $A$. The following statements are equivalent:
    (i) $A$ is primitive, (ii) $G$ is strongly connected and aperiodic.
  \end{block}
  \begin{block}{Corollary 1.12 (Primitive matrices and digraphs with self-loops) [2]}
    Let $G = (V, E, \{a_{ij}\}_{(i,j) \in E})$ be a weighted digraph of order $n \geq 2$
    with adjacency matrix $A$. If $G$ has at least one self-loop and is strongly connected,
    then $A$ is primitive.
  \end{block}
  \detailbuttons{\go{bk-prelim3}{Next}\ \ret{m-fdg-model}}
\end{frame}

\begin{frame}[label=bk-prelim3]{Row-substochastic matrices}
  \scriptsize
  \begin{block}{Definition 1.13 (Row-substochastic matrix) [1]}
    A non-negative matrix $A \in \mathbb{R}^{n \times n}$ is \emph{row-substochastic} if
    its row-sums are at most $1$ and at least one row-sum is strictly less than $1$. That
    is: $A \mathbf{1}_n \leq \mathbf{1}_n$ and
    $\exists\, i \in \{1, \dots, n\}$ s.t.\ $e_i^\top A \mathbf{1}_n < 1$.
  \end{block}
  \begin{block}{Lemma 1.14 (Convergent row-substochastic matrices) [1]}
    Let $G$ be a weighted digraph with row-substochastic adjacency matrix $A$. Then, $A$
    is convergent if and only if $G$ contains directed walks from each node with
    (weighted) out-degree $1$ to a node with (weighted) out-degree strictly less than $1$.
  \end{block}
  \detailbuttons{\ret[Back to the FDG model]{m-fdg-model}\ \ret[Back to FJ stability]{bk-fj-stab}}
\end{frame}

% ---------------------------------------------------------------------------
% FDG model
% ---------------------------------------------------------------------------

\begin{frame}[label=bk-fdg-thm]{Theorem 2.2 (consensus)}
  \scriptsize
  \begin{block}{Definition 2.1 (Convergence) [5]}
    The FDG model (1) is convergent if for any initial condition $x(0)$ the limit exists
    \[ x(\infty) = \lim_{k \to \infty} x(k) = \lim_{k \to \infty} A^k x(0). \]
    A convergent model reaches a consensus if $x_1(\infty) = \dots = x_n(\infty)$ for any
    initial opinion vector $x(0)$.
  \end{block}
  \begin{block}{Theorem 2.2 (Consensus in DT averaging systems) [1]}
    Let $G$ be a weighted digraph with row-stochastic adjacency matrix $A$. The following
    statements are equivalent:
    \begin{itemize}
      \item[(A1)] the eigenvalue $1$ is simple and all other eigenvalues $\mu$ satisfy
                  $|\mu| < 1$;
      \item[(A2)] $A$ is semi-convergent and $\lim_{k \to \infty} A^k = \mathbf{1}_n w^\top$
                  where $w \in \mathbb{R}^n$ satisfies $w \geq 0$,
                  $\mathbf{1}_n^\top w = 1$, and $w^\top A = w^\top$;
      \item[(A3)] $G$ contains a globally reachable node and the subgraph of globally
                  reachable nodes is aperiodic.
    \end{itemize}
    If any, and therefore all, of the previous conditions are satisfied, then
    \begin{itemize}
      \item[(i)]  $w$ is the left dominant eigenvector of $A$ and $w_i > 0$ iff node $i$ is
                  globally reachable;
      \item[(ii)] the solution of the DT averaging system satisfies
                  $\lim_{k \to \infty} x(k) = \big(w^\top x(0)\big)\mathbf{1}_n$.
    \end{itemize}
  \end{block}
  \detailbuttons{\go{bk-fdg-thm2}{Disagreement}\ \ret[Back to the FDG model]{m-fdg-model}\ \ret[Back to social power]{m-fdg-power}}
\end{frame}

\begin{frame}[label=bk-fdg-thm2]{Average consensus, disagreement}
  \scriptsize
  \begin{block}{Corollary 2.3 (Average consensus) [1]}
    If $A$ is doubly-stochastic and (A3) holds, then $w = \frac{1}{n}\mathbf{1}_n$ and
    $\lim_{k \to +\infty} x(k) = \big(\frac{1}{n}\sum_{i=1}^{n} x_i(0)\big)\, \mathbf{1}_n$.
  \end{block}
  \begin{block}{Theorem 2.4 (Disagreement in DT averaging systems) [1]}
    Let $G$ be a weighted digraph with row-stochastic adjacency matrix $A$. Assume that the
    condensation digraph $C(G)$ has $n_s \geq 2$ sinks. The following statements are
    equivalent:
    \begin{itemize}
      \item[(A1)] the eigenvalue $1$ is semi-simple with multiplicity $n_s$ and all other
                  eigenvalues $\mu$ satisfy $|\mu| < 1$;
      \item[(A2)] $A$ is semi-convergent;
      \item[(A3)] each sink of $C(G)$ is aperiodic.
    \end{itemize}
    If any, and therefore all, of the previous conditions are satisfied, then the solution
    of the DT averaging system satisfies
    $\lim_{k \to \infty} x(k) = v \neq \alpha \mathbf{1}_n$, $\alpha \in \mathbb{R}$.
  \end{block}
  \textbf{Group by group [1].} Each sink evolves as an FDG model on its own and, being
  aperiodic, reaches an internal consensus that depends only on the initial opinions of its
  members; every agent outside the sinks converges to a convex combination of the consensus
  values of the sinks.
  \detailbuttons{\ret[Back to Theorem 2.2]{bk-fdg-thm}\ \ret[Back to the FDG model]{m-fdg-model}}
\end{frame}

\begin{frame}[label=bk-rosetta]{Rosetta stone (Table 1)}
  \centering
  \scriptsize
  \renewcommand{\arraystretch}{1.3}
  \begin{tabular}{@{}c c c c >{\raggedright\arraybackslash}m{4.6cm} l@{}}
    \toprule
    Case & Example digraph & $A$ & $\operatorname{spec}(A)$ & Property of $A$ / of $G$ & Outcome \\
    \midrule
    (i)
      & \begin{tikzpicture}[baseline=-0.6ex]
          \node[agent] (a) at (0,0) {1}; \node[agent] (b) at (1.4,0) {2};
          \draw[edge] (a) to[bend left=25] node[weight, above] {1/2} (b);
          \draw[edge] (b) to[bend left=25] node[weight, below] {1/2} (a);
          \draw[selfloop] (a) to node[weight, above] {1/2} (a);
          \draw[selfloop] (b) to node[weight, above] {1/2} (b);
        \end{tikzpicture}
      & $\left[\begin{smallmatrix} 1/2 & 1/2 \\ 1/2 & 1/2 \end{smallmatrix}\right]$
      & $\{1,\, 0\}$
      & $A$ primitive / $G$ strongly connected and aperiodic
      & consensus \\
    (ii)
      & \begin{tikzpicture}[baseline=-0.6ex]
          \node[agent] (a) at (0,0) {1}; \node[agent] (b) at (1.4,0) {2};
          \draw[edge] (a) to[bend left=25] (b);
          \draw[edge] (b) to[bend left=25] (a);
        \end{tikzpicture}
      & $\left[\begin{smallmatrix} 0 & 1 \\ 1 & 0 \end{smallmatrix}\right]$
      & $\{1,\, -1\}$
      & $A$ irreducible, not primitive / $G$ strongly connected and periodic
      & no convergence \\
    (iii)
      & \begin{tikzpicture}[baseline=-0.6ex]
          \node[agent] (a) at (0,0) {1}; \node[agent] (b) at (1.4,0) {2};
          \draw[edge] (b) to node[weight, below] {1/2} (a);
          \draw[selfloop] (a) to (a);
          \draw[selfloop] (b) to node[weight, above] {1/2} (b);
        \end{tikzpicture}
      & $\left[\begin{smallmatrix} 1 & 0 \\ 1/2 & 1/2 \end{smallmatrix}\right]$
      & $\{1,\, 1/2\}$
      & $1$ simple and strictly dominant, $A$ not primitive / $C(G)$ has a unique sink, which is aperiodic
      & consensus \\
    (iv)
      & \begin{tikzpicture}[baseline=-0.6ex]
          \node[agent] (a) at (0,0) {1}; \node[agent] (b) at (1.4,0) {2};
          \draw[selfloop] (a) to (a);
          \draw[selfloop] (b) to (b);
        \end{tikzpicture}
      & $\left[\begin{smallmatrix} 1 & 0 \\ 0 & 1 \end{smallmatrix}\right]$
      & $\{1,\, 1\}$
      & $1$ semi-simple but not simple, strictly dominant / $C(G)$ has at least two sinks, all aperiodic
      & disagreement \\
    (v)
      & \begin{tikzpicture}[baseline=-0.6ex]
          \node[agent] (a) at (0,0) {1}; \node[agent] (b) at (1.1,0) {2};
          \node[agent] (c) at (2.2,0) {3};
          \draw[edge] (a) to[bend left=25] (b);
          \draw[edge] (b) to[bend left=25] (a);
          \draw[selfloop] (c) to (c);
        \end{tikzpicture}
      & $\left[\begin{smallmatrix} 0 & 1 & 0 \\ 1 & 0 & 0 \\ 0 & 0 & 1 \end{smallmatrix}\right]$
      & $\{1,\, 1,\, -1\}$
      & $1$ not strictly dominant / $G$ not strongly connected, at least one sink of $C(G)$ is periodic
      & no convergence \\
    \bottomrule
  \end{tabular}

  \vspace{1mm}
  {\tiny Cases (i)--(iv) and the examples follow [1]; the properties of case (v), left
  blank in the slides, follow from [5].}
  \detailbuttons{\ret{m-fdg-model}}
\end{frame}

\begin{frame}[label=bk-stubborn]{Stubborn agents in the FDG}
  \scriptsize
  \begin{block}{Definition 2.5 (Stubbornness) [5]}
    An agent is said to be \emph{stubborn} if its opinion remains unchanged independent of
    the others' opinions.
  \end{block}
  If the opinions obey the FDG model, then agent $i$ is stubborn, $x_i(k) \equiv x_i(0)$,
  if and only if $a_{ii} = 1$ [5] (and hence $a_{ij} = 0$ for all $j \neq i$, since $A$ is
  row-stochastic). In the digraph, a stubborn agent is a strongly connected component made
  of a single node, which is a sink of $C(G)$ and is aperiodic because of its self-loop.
  If every agent is influenced by the stubborn agent $i$, then $i$ is the only globally
  reachable node, and Theorem~2.2 gives $w = e_i$.
  \begin{block}{Corollary 2.6 [5]}
    Let the group have $s \geq 1$ stubborn agents, influencing all other individuals (i.e.\
    every other node of $G$ is connected by a walk to the set of stubborn nodes). Then the
    FDG model (1) is convergent.
  \end{block}
  \begin{block}{Remark 2.7 (More than one stubborn agent)}
    If two or more agents are stubborn, each of them is a sink of $C(G)$ made of a single,
    aperiodic node. If every other agent is influenced, directly or through intermediaries,
    by at least one stubborn agent, these are the only sinks, and the group converges
    (Corollary~2.6) to a disagreement (Theorem~2.4): every stubborn agent keeps its own
    initial opinion, and every other agent converges to a convex combination of the
    opinions of the stubborn agents it is influenced by. Since there is no single consensus
    value, the social power is no longer defined.
  \end{block}
  \detailbuttons{\ret{m-fdg-power}}
\end{frame}

\begin{frame}[label=bk-spectra]{Spectra of $A(t)$}
  \small
  \begin{itemize}
    \item $t = 0$. Directed cycle: permutation matrix with $A^6 = I$, so its spectrum is the
          set of the sixth roots of unity. Bidirectional ring: circulant, with eigenvalues
          $\cos(2\pi k/6)$, $k = 0, \dots, 5$, i.e.\
          $\{1, \frac12, \frac12, -\frac12, -\frac12, -1\}$. Some eigenvalues different from
          $1$ lie on the unit circle: they cause the oscillations.
    \item $t > 0$: all the eigenvalues different from $1$ lie strictly inside the unit
          disk, and the group converges. For $t = 1/2$ the largest of their moduli,
          computed numerically, is $\approx 0.92$ on the directed cycle, where the
          characteristic polynomial of $A(1/2)$ is
          $\lambda^5(\lambda - \tfrac12) - \tfrac12$, and $\approx 0.90$ on the
          bidirectional ring.
    \item $t = 1$, directed cycle, agent~$1$ listed first: $A(1)$ is block lower
          triangular,
          \[ A(1) = \begin{bmatrix} 1 & 0 \\ b & N \end{bmatrix}, \qquad
             b = (0, 0, 0, 0, 1)^\top, \]
          where $b$ records that agent~$6$ listens to agent~$1$ and the block $N$ of agents
          $2, \dots, 6$ is strictly upper triangular, since agent~$i$ copies agent~$i+1$.
          Hence $N^5 = 0$, the eigenvalue $0$ has multiplicity $5$, and
          $A(1)^5 = \mathbf{1}_6 e_1^\top$: after five steps every agent holds $x_1(0)$.
  \end{itemize}
  \detailbuttons{\ret{m-fdg-conv}}
\end{frame}

\begin{frame}[label=bk-socialpower]{Social power $w(t)$}
  \small
  \begin{itemize}
    \item $0 < t < 1$: strongly connected, every agent globally reachable, so every agent
          has a positive social power (Theorem~2.2(i)). $t = 1$: agent~$1$ is the only
          globally reachable node, $w = e_1$.
    \item In between, $w$ is the left eigenvector of $A(t)$ for the eigenvalue $1$,
          normalized to sum $1$ (condition (A2) of Theorem~2.2). Substituting
          $w = (w_1, b, \dots, b)$ in $w^\top A(t) = w^\top$: in both networks agents
          $2, \dots, 6$ receive the same total weight from the others, and the equations
          reduce to $b = (1-t)\, w_1$. Hence
          \[ w(t) = \frac{1}{1 + 5(1-t)}\,\big(1,\ 1-t,\ 1-t,\ 1-t,\ 1-t,\ 1-t\big)^\top . \]
    \item $w_1 = 1/(1 + 5(1-t))$ grows from $1/6$ to $1$; for $t = 1/2$, $w_1 = 2/7$ and
          $w_i = 1/7$ for $i \neq 1$.
    \item Consensus value (Theorem~2.2(ii)), with $x(0) = (0, 0.2, 0.4, 0.6, 0.8, 1)^\top$:
          \[ w(t)^\top x(0) = \frac{3(1-t)}{6 - 5t}, \]
          from the average $0.5$ for $t \to 0$ to $x_1(0) = 0$ for $t = 1$, and $3/7$ for
          $t = 1/2$.
  \end{itemize}
  \detailbuttons{\ret{m-fdg-sp}}
\end{frame}

% ---------------------------------------------------------------------------
% FJ model: theory
% ---------------------------------------------------------------------------

\begin{frame}[label=bk-fj-thm]{Theorem 3.2}
  \scriptsize
  \begin{block}{Definition 3.1 (P-dependent agents) [5]}
    Agent $i$ is \emph{P-dependent} if it is prejudiced or influenced by some prejudiced
    agent $j$, that is, a walk from $i$ to $j$ in the graph $G$ exists. Otherwise, the agent
    is \emph{P-independent}.
  \end{block}
  Agents $1, \dots, r$ P-dependent, $r+1, \dots, n$ P-independent: $W^{21} = 0$,
  $\Lambda^{22} = I$, and [5]
  \[
    x^1(k+1) = \Lambda^{11}\bigl[W^{11} x^1(k) + W^{12} x^2(k)\bigr]
               + (I - \Lambda^{11})\, u^1 \;\; (6), \qquad
    x^2(k+1) = W^{22} x^2(k) \;\; (7).
  \]
  \begin{block}{Theorem 3.2 [5]}
    Let the community have $r \geq 1$ P-dependent agents and $n - r \geq 0$ P-independent
    ones. Then, the subsystem (6) is asymptotically stable, i.e.\ $\Lambda^{11} W^{11}$ is
    Schur stable $\rho(\Lambda^{11} W^{11}) < 1$. The model (4) is convergent if and only if
    $r = n$ or (7) is convergent, i.e.\ $W^{22}$ is regular. In this case
    \[
      x^1(\infty) = V \begin{bmatrix} u^1 \\ x^2(\infty) \end{bmatrix}, \qquad
      V \triangleq (I_r - \Lambda^{11} W^{11})^{-1}
          \begin{bmatrix} I_r - \Lambda^{11} & \Lambda^{11} W^{12} \end{bmatrix}.
    \]
    The matrix $V$ is stochastic, i.e.\ the final opinion of any agent is a convex
    combination of the prejudices and final opinions of P-independent agents.
  \end{block}
  \detailbuttons{\go{bk-fj-stab}{Stability}\ \ret{m-fj-conv}}
\end{frame}

\begin{frame}[label=bk-fj-stab]{Asymptotic stability}
  \scriptsize
  \begin{block}{Corollary 3.3 [5]}
    The Friedkin--Johnsen model (4) is asymptotically stable if and only if all agents are
    P-dependent. Then, $V = (I - \Lambda W)^{-1}(I - \Lambda)$. This holds, in particular,
    if $\Lambda < I_n$ or $\Lambda \neq I_n$ and $G$ is strongly connected ($W$ is
    irreducible).
  \end{block}
  \textbf{With the tools of the course.} Row $i$ of $\Lambda W$ sums to $\lambda_i$: rows of
  agents that are not prejudiced sum to $1$, rows of prejudiced agents to less than $1$, so
  $\Lambda W$ is row-substochastic (Definition~1.13) and Lemma~1.14 applies: every agent that
  is not prejudiced must be influenced, directly or through intermediaries, by a prejudiced
  agent, i.e.\ all agents are P-dependent.

  \medskip
  \textbf{Intuition.} $x(\infty) = Vu$ satisfies $x(\infty) = \Lambda W x(\infty) + (I -
  \Lambda)u$, since $(I - \Lambda W)V = I - \Lambda$. With $e(k) = x(k) - x(\infty)$:
  \[ e(k+1) = \Lambda W\, e(k) \qquad (10), \]
  an autonomous system without $u$: $\Lambda W$ decides whether the opinions converge, $u$
  decides where. Componentwise $e_i(k+1) = \lambda_i \sum_j w_{ij}\, e_j(k)$: a prejudiced
  agent loses a fraction $1 - \lambda_i$ of the deviations, so they vanish only if every
  agent is influenced by a prejudiced agent.
  \detailbuttons{\go{bk-fj-stab2}{Convergence, spectrum}\ \go{bk-prelim3}{Lemma 1.14}\ \ret{m-fj-conv}}
\end{frame}

\begin{frame}[label=bk-fj-stab2]{Convergence and spectrum}
  \footnotesize
  \begin{block}{Corollary 3.4 [5]}
    The model (4) converges if and only if $\Lambda W$ is regular, i.e.\ the limit
    $\lim_{k \to \infty} (\Lambda W)^k$ exists.
  \end{block}
  \textbf{Theorem 3.2 by $r$.} $r = 0$: not covered, $\Lambda = I$ is the FDG model.
  $0 < r < n$: the P-independent agents follow the FDG model with $W^{22}$, read on the
  sinks of $C^{22}$. $r = n$: no P-independent agents, the model always converges
  (Corollary~3.3).

  \medskip
  Not trivial [5]: a system $x(k+1) = A x(k) + B u$ with a regular matrix $A$ may not
  converge, and may even have unbounded solutions, as for $A = B = I$.

  \medskip
  \textbf{Spectrum.} Listing the P-dependent agents first,
  \[
    \Lambda W = \begin{bmatrix} \Lambda^{11} W^{11} & \Lambda^{11} W^{12} \\ 0 & W^{22}
    \end{bmatrix},
    \qquad
    \operatorname{spec}(\Lambda W) = \operatorname{spec}(\Lambda^{11} W^{11}) \cup
    \operatorname{spec}(W^{22}),
  \]
  and $\rho(\Lambda^{11} W^{11}) < 1$ (Theorem~3.2): the eigenvalues on the unit circle
  can only come from $W^{22}$, i.e.\ from the FDG model of the P-independent agents.
  Asymptotically stable when $r = n$; convergent when the eigenvalue $1$ of $W^{22}$ is
  semi-simple and its other eigenvalues lie inside the unit disk, i.e.\ $W^{22}$ regular.
  \detailbuttons{\ret[Back to stability]{bk-fj-stab}\ \ret{m-fj-conv}}
\end{frame}

\begin{frame}[label=bk-fj-consensus]{Consensus in the FJ model}
  \small
  All agents P-dependent: $x_i(\infty) = \sum_{j=1}^{n} v_{ij}\, u_j$, so
  \begin{itemize}
    \item the initial opinions are forgotten: $x(0)$ does not appear in $x(\infty)$;
    \item only the prejudices of prejudiced agents count: if $\lambda_j = 1$, the $j$-th
          column of $I - \Lambda$ is zero, hence so is the $j$-th column of $V$;
    \item every final opinion lies between the smallest and the largest prejudice of the
          prejudiced agents.
  \end{itemize}
  For a given $u$, consensus if $x(\infty) = \gamma\,\mathbf{1}_n$. Since
  $W\mathbf{1}_n = \mathbf{1}_n$, $\gamma\mathbf{1}_n$ is an equilibrium iff
  \[ \gamma\,\mathbf{1}_n = \Lambda W \gamma\,\mathbf{1}_n + (I - \Lambda)\, u
     \iff (I - \Lambda)(\gamma\,\mathbf{1}_n - u) = 0, \]
  i.e.\ $(1 - \lambda_i)(\gamma - u_i) = 0$ for every agent $i$; when all agents are
  P-dependent the equilibrium is unique and equal to $Vu$ (Corollary~3.3).
  \begin{alertblock}{}
    When all agents are P-dependent, the FJ model reaches a consensus if and only if all
    the prejudiced agents have the same prejudice, and the group agrees on it.
  \end{alertblock}
  \detailbuttons{\ret{m-fj-final}}
\end{frame}

\begin{frame}[label=bk-augmented]{The augmented model}
  \small
  Each agent $i$ gets a virtual agent $i'$ with $x_{i'}(k) = u_i$ for all $k$ [5]:
  \[
    \begin{bmatrix} x(k+1) \\ x'(k+1) \end{bmatrix}
    = \underbrace{\begin{bmatrix} \Lambda W & I - \Lambda \\ 0 & I \end{bmatrix}}_{\bar A}
      \begin{bmatrix} x(k) \\ x'(k) \end{bmatrix},
    \qquad x'(0) = u . \qquad (11)
  \]
  Row $i$ of $\bar A$ sums to $\lambda_i + (1 - \lambda_i) = 1$, row $i'$ to $1$: an FDG
  model on $2n$ agents. In the augmented graph a real agent $i$ listens to $j$ with weight
  $\lambda_i w_{ij}$ and to $i'$ with weight $1 - \lambda_i$; $i'$ listens only to itself:
  it is stubborn.
  \begin{itemize}
    \item Every virtual agent is an aperiodic sink: at least $n$ sinks, so cases (i),
          (ii), (iii) never occur.
    \item \emph{Disagreement} (case (iv)) when all sinks are aperiodic: with all agents
          P-dependent the only sinks are the virtual agents, and every real agent converges
          to a convex combination of the prejudices, i.e.\ (9); otherwise the closed groups
          of P-independent agents are further sinks, all aperiodic iff $W^{22}$ is regular.
    \item \emph{No convergence} (case (v)) when a sink of $C^{22}$ is periodic.
  \end{itemize}
  \detailbuttons{\go{bk-augmented2}{Two descriptions}\ \ret{m-fj-final}}
\end{frame}

\begin{frame}[label=bk-augmented2]{Two descriptions, one dynamics}
  \small
  With all agents P-dependent, the FJ model is asymptotically stable (Corollary~3.3),
  while the augmented model (11), row-stochastic, is never asymptotically stable: it
  converges to the disagreement of case (iv), which depends on $x'(0) = u$. $\bar A$ is
  block triangular, so
  \[
    \operatorname{spec}(\bar A) = \operatorname{spec}(\Lambda W) \cup
    \{\underbrace{1, \dots, 1}_{n}\}.
  \]
  \begin{itemize}
    \item The eigenvalues of $\Lambda W$ lie inside the unit disk: the FJ model, which has
          only these, is asymptotically stable.
    \item The augmented model adds the eigenvalue $1$ once for each virtual agent,
          semi-simple by Theorem~2.4 applied to its $n$ aperiodic sinks: the prejudices,
          which in the augmented model are \alert{states}.
    \item In the FJ model $u$ is an \alert{input}: asymptotic stability concerns the
          perturbations of $x(0)$ with $u$ fixed, which die out as in (10); changing $u$
          is a different input, with a different equilibrium $Vu$.
  \end{itemize}
  \detailbuttons{\ret[Back to the augmented model]{bk-augmented}\ \ret{m-fj-final}}
\end{frame}

\begin{frame}[label=bk-centrality]{Influence centrality}
  \small
  Asymptotically stable FJ model with $u = x(0)$ [5]: $x(\infty) = V x(0)$, and
  \[ c = n^{-1} V^\top \mathbf{1}_n, \qquad
     \bar x = \frac{1}{n} \sum_{i=1}^{n} x_i(\infty) = \frac{1}{n} \mathbf{1}_n^\top V x(0)
            = c^\top x(0). \qquad (12) \]
  $c_i$ is the average of the $i$-th column of $V$; $V$ stochastic gives $c \geq 0$ and
  $c^\top \mathbf{1}_n = 1$. For $\Lambda = \alpha I_n$, $\alpha \in [0,1)$: for $\alpha = 0$,
  $V = I$ and $c_0 = n^{-1}\mathbf{1}_n$ [5].
  \begin{block}{Lemma 3.5 [5]}
    Let $W$ be a fully regular matrix and $w$ stand for the vector of French's social power,
    corresponding to the FDG model (1) with $A = W$. Then, $w = \lim_{\alpha \to 1-0}
    c_\alpha$.
  \end{block}
  As $\alpha \to 1$, $V_\alpha = (1 - \alpha)(I - \alpha W)^{-1} \to \mathbf{1}_n w^\top$
  [5]: the final opinions approach the FDG consensus, and $c_\alpha \to w$. The limit
  cannot be replaced by $\alpha = 1$: for $\Lambda = I$ the FJ model is the FDG model, which
  is not asymptotically stable, and $I - W$ is singular since $1 \in \operatorname{spec}(W)$
  (Theorem~1.6), so $V$ is not defined.
  \detailbuttons{\ret{m-fj-centrality}}
\end{frame}

% ---------------------------------------------------------------------------
% FJ model: examples
% ---------------------------------------------------------------------------

\begin{frame}[label=bk-state-input]{How a perturbation evolves}
  \small
  $x(k)$, $\tilde x(k)$: trajectories of Figures 6 and 7, $\delta(k) = \tilde x(k) - x(k)$.
  \begin{itemize}
    \item FJ: $u$ is the same in the two runs, so $\delta(k+1) = \Lambda W\, \delta(k)$, the
          recursion of the error (10). All agents P-dependent: $\Lambda W$ Schur stable,
          $\delta(k) \to 0$, the two runs reach the same $Vu$.
    \item FDG: $\delta(k+1) = A\,\delta(k)$ and, by Theorem~2.2,
          $\delta(k) \to (w^\top \delta(0))\mathbf{1}_6 = 0.5\,\mathbf{1}_6$, since
          $w = e_1$ and $\delta(0) = 0.5\, e_1$: never forgotten.
  \end{itemize}
  \textbf{Spectra.} $\Lambda W$ and $A(1)$ differ only in the row of agent~$1$: zero in
  $\Lambda W$, $e_1^\top$ in $A(1)$. Listing agent~$1$ first, both are block lower
  triangular, with diagonal entry $0$ or $1$ for agent~$1$ and the same block for agents
  $2, \dots, 6$.
  \begin{itemize}
    \item Ring: the block is $\tfrac12$ times the adjacency matrix of the path graph on
          five nodes, whose spectrum is $\{2\cos(k\pi/6) \mid k = 1, \dots, 5\}$ [4], so the
          block has eigenvalues $\cos(k\pi/6)$:
          $\operatorname{spec}(\Lambda W) = \{0, 0, \pm\tfrac12, \pm\tfrac{\sqrt3}{2}\}$,
          $\operatorname{spec}(A(1)) = \{1, 0, \pm\tfrac12, \pm\tfrac{\sqrt3}{2}\}$.
    \item Cycle: the block is the nilpotent $N$: all eigenvalues $0$, except the eigenvalue
          $1$ of $A(1)$, that of the virtual agent $1'$.
  \end{itemize}
  \detailbuttons{\ret{m-state-input}}
\end{frame}

\begin{frame}[label=bk-rewritings]{From one model to the other}
  \centering
  \begin{tikzpicture}[baseline=(c.base), scale=0.75]
    \foreach \i in {1,...,6} {
      \node[agent] (rd\i) at ({90-60*(\i-1)}:1.3) {\i};
    }
    \node (c) at (0,0) {};
    \foreach \i/\j in {2/3,3/4,4/5,5/6,6/1} {
      \draw[edge] (rd\i) -- (rd\j);
    }
    \draw[selfloop] (rd1) to node[weight, above] {$1$} (rd1);
    \node at (0,-2.0) {\scriptsize (a) FDG model, $a_{11} = 1$};
  \end{tikzpicture}
  \hspace{1.5cm}
  \begin{tikzpicture}[baseline=(c.base), scale=0.75]
    \foreach \i in {1,...,6} {
      \node[agent] (rf\i) at ({90-60*(\i-1)}:1.3) {\i};
    }
    \node (c) at (0,0) {};
    \foreach \i/\j in {2/3,3/4,4/5,5/6,6/1} {
      \draw[edge] (rf\i) -- (rf\j);
    }
    \draw[edge, dashed, gray] (rf1) -- node[weight, above right] {$\lambda_1 = 0$} (rf2);
    \draw[edge] (0,2.7) node[above, font=\scriptsize] {$u_1 = x_1(0)$}
      -- node[weight, right] {$1-\lambda_1 = 1$} (rf1);
    \node at (0,-2.0) {\scriptsize (b) FJ model, $\lambda_1 = 0$};
  \end{tikzpicture}

  {\scriptsize Figure 5: the rewriting of [5] on the directed cycle with $t = 1$.}

  \vspace{2mm}
  \raggedright
  \footnotesize
  \begin{itemize}
    \item $\lambda_i = 0$ and $u_i = x_i(0)$ for every stubborn agent, $\lambda_i = 1$ (and
          any $u_i$) for the others: the two models produce the same trajectories.
    \item The row of $W$ of a stubborn agent plays no role, since it is multiplied by
          $\lambda_i = 0$. The stubborn agent turns from a state into an input; its opinion
          appears twice, as $x_i(0)$ and as $u_i = x_i(0)$.
  \end{itemize}
  \detailbuttons{\go{bk-rewritings2}{Second rewriting}\ \ret{m-state-input}}
\end{frame}

\begin{frame}[label=bk-rewritings2]{The stubborn agent as a prejudice}
  \centering
  \begin{tikzpicture}[baseline=(c.base), scale=0.75]
    \foreach \i in {1,...,6} {
      \node[agent] (sa\i) at ({90-60*(\i-1)}:1.3) {\i};
    }
    \node (c) at (0,0) {};
    \foreach \i/\j in {2/3,3/4,4/5,5/6} {
      \draw[edge] (sa\i) to[bend left=12] (sa\j);
      \draw[edge] (sa\j) to[bend left=12] (sa\i);
    }
    \draw[edge] (sa2) -- (sa1);
    \draw[edge] (sa6) -- (sa1);
    \draw[selfloop] (sa1) to node[weight, above] {$1$} (sa1);
    \node at (0,-2.0) {\scriptsize (a) FDG model, agent~$1$ stubborn};
  \end{tikzpicture}
  \hspace{1.5cm}
  \begin{tikzpicture}[baseline=(c.base), scale=0.75]
    \foreach \i in {2,...,6} {
      \node[agent] (sb\i) at ({90-60*(\i-1)}:1.3) {\i};
    }
    \node (c) at (0,0) {};
    \node[font=\scriptsize] (su) at (90:1.3) {$x_1(0)$};
    \foreach \i/\j in {3/4,4/5} {
      \draw[edge] (sb\i) to[bend left=12] (sb\j);
      \draw[edge] (sb\j) to[bend left=12] (sb\i);
    }
    \draw[edge] (sb2) to[bend left=12] node[weight, right] {$\lambda_2$} (sb3);
    \draw[edge] (sb3) to[bend left=12] (sb2);
    \draw[edge] (sb6) to[bend right=12] node[weight, left] {$\lambda_6$} (sb5);
    \draw[edge] (sb5) to[bend right=12] (sb6);
    \draw[edge] (su) -- node[weight, above right] {$1-\lambda_2$} (sb2);
    \draw[edge] (su) -- node[weight, above left] {$1-\lambda_6$} (sb6);
    \node at (0,-2.0) {\scriptsize (b) FJ model on agents $2, \dots, 6$};
  \end{tikzpicture}

  {\scriptsize Figure 8: bidirectional ring with $t = 1$; every edge between agents has
  weight $1/2$.}

  \vspace{2mm}
  \raggedright
  \footnotesize
  \begin{itemize}
    \item An agent $i$ that listens to the stubborn agent~$1$ updates as
          $x_i(k+1) = \sum_{j \neq 1} a_{ij}\, x_j(k) + a_{i1}\, x_1(0)$, which is (5) with
          $\lambda_i = 1 - a_{i1}$, $w_{ij} = a_{ij}/\lambda_i$, $u_i = x_1(0)$.
    \item On the ring: $\lambda_2 = \lambda_6 = 1/2$, $u_2 = u_6 = x_1(0)$,
          $\lambda_3 = \lambda_4 = \lambda_5 = 1$: the augmented graph read backwards. The
          stubborn opinion appears only once, as a prejudice.
  \end{itemize}
  \detailbuttons{\ret[Back to the first rewriting]{bk-rewritings}\ \ret{m-state-input}}
\end{frame}

\begin{frame}[label=bk-one-prej]{A single prejudice}
  \small
  On the directed cycle:
  \[
    {\setlength{\arraycolsep}{3pt}
    \Lambda W = \begin{bmatrix}
       0&\frac12&0&0&0&0\\ 0&0&1&0&0&0\\ 0&0&0&1&0&0\\ 0&0&0&0&1&0\\ 0&0&0&0&0&1\\ 1&0&0&0&0&0
    \end{bmatrix},
    \qquad
    A(\tfrac12) = \begin{bmatrix}
       \frac12&\frac12&0&0&0&0\\ 0&0&1&0&0&0\\ 0&0&0&1&0&0\\ 0&0&0&0&1&0\\ 0&0&0&0&0&1\\ 1&0&0&0&0&0
    \end{bmatrix}}
  \]
  The weight $1/2$ taken off the agent it listens to goes, in $A(1/2)$, to the diagonal
  (current opinion), and in the FJ model out of the matrix, to the input
  $(1 - \lambda_1)u_1$; the first row of $\Lambda W$ sums to $1/2$.
  \begin{itemize}
    \item Strongly connected: all agents P-dependent, asymptotically stable
          (Corollary~3.3, case $\Lambda \neq I_n$, $W$ irreducible).
    \item Column $j$ of $V$ is zero when $\lambda_j = 1$ and $V$ is stochastic:
          $V = \mathbf{1}_6 e_1^\top$, $x(\infty) = u_1\mathbf{1}_6 = 0$, $c = e_1$.
    \item Augmented graph: $2', \dots, 6'$ have weight $0$, nobody listens to them; the
          rest has the single sink $1'$, so it reaches a consensus with social power
          $e_{1'}$: with $u = x(0)$, the group converges to $x_1(0)$.
  \end{itemize}
  \detailbuttons{\ret{m-one-prej}}
\end{frame}

\begin{frame}[label=bk-all-prej]{Every agent prejudiced}
  \small
  On the directed cycle:
  \[
    {\setlength{\arraycolsep}{3pt}
    A = 0.2\, I + 0.8\, W = \begin{bmatrix}
       0.2&0.8&0&0&0&0\\ 0&0.2&0.8&0&0&0\\ 0&0&0.2&0.8&0&0\\ 0&0&0&0.2&0.8&0\\
       0&0&0&0&0.2&0.8\\ 0.8&0&0&0&0&0.2
    \end{bmatrix},
    \qquad
    \Lambda W = 0.8\, W = \begin{bmatrix}
       0&0.8&0&0&0&0\\ 0&0&0.8&0&0&0\\ 0&0&0&0.8&0&0\\ 0&0&0&0&0.8&0\\
       0&0&0&0&0&0.8\\ 0.8&0&0&0&0&0
    \end{bmatrix}}
  \]
  \begin{itemize}
    \item FDG: $A$ row- and column-stochastic with positive diagonal, primitive
          (Corollary~1.12): consensus on the average $0.5$ (Corollary~2.3).
    \item FJ: $\Lambda < I_6$, asymptotically stable (Corollary~3.3),
          $x(\infty) = Vu$ with $V = 0.2\,(I - 0.8\,W)^{-1}$, stochastic, with all entries
          positive since $W$ is irreducible: every prejudice enters every final opinion,
          with weights that change from row to row; the prejudices differ, so no consensus.
    \item Augmented graph: six stubborn virtual agents with different opinions, as in
          Remark~2.7.
  \end{itemize}
  \detailbuttons{\ret{m-all-prej}}
\end{frame}

\begin{frame}[label=bk-aux]{Auxiliary network: blocks}
  \scriptsize
  Order $\sigma = (4, 5 \mid 1, 2, 3, 6)$, $P = [\, e_4\ e_5\ e_1\ e_2\ e_3\ e_6 \,]^\top$; for
  $W_{\mathrm{a}}$:
  \[
    \tilde W_{\mathrm{a}} = P\, W_{\mathrm{a}} P^\top =
    \left[\begin{array}{cc|cccc}
      0 & 1 & 0 & 0 & 0 & 0\\
      0 & 0 & 0 & 0 & 0 & 1\\ \hline
      0 & 0 & 0 & 1 & 0 & 0\\
      0 & 0 & \tfrac12 & 0 & \tfrac12 & 0\\
      0 & 0 & 1 & 0 & 0 & 0\\
      0 & 0 & 1 & 0 & 0 & 0
    \end{array}\right]
    = \begin{bmatrix} W^{11} & W^{12} \\ 0 & W^{22} \end{bmatrix},
    \qquad
    \tilde\Lambda = \mathrm{diag}(1,\, 0.5 \mid 1,\, 1,\, 1,\, 1).
  \]
  The lower-left block is zero: no P-independent agent listens to a P-dependent one. With
  $\Lambda^{11} = \mathrm{diag}(1, 0.5)$, Theorem~3.2 gives
  \[
    V = (I_2 - \Lambda^{11} W^{11})^{-1}\big[\, I_2 - \Lambda^{11},\ \Lambda^{11} W^{12} \,\big]
      = \left[\begin{array}{cc|cccc}
          0 & 0.5 & 0 & 0 & 0 & 0.5\\ 0 & 0.5 & 0 & 0 & 0 & 0.5
        \end{array}\right],
    \quad x_4(\infty) = x_5(\infty) = 0.5\, u_5 + 0.5\, x_6(\infty).
  \]
  Social power of the closed group $\{1, 2, 3\}$, from its block of $W^{22}$
  (Theorem~2.2): $w_1 = \tfrac12 w_2 + w_3$, $w_2 = w_1$, $w_3 = \tfrac12 w_2$, sum $1$
  $\Rightarrow$ $(0.4, 0.4, 0.2)$.
  \begin{itemize}
    \item (c): consensus $0.4 \cdot 0 + 0.4 \cdot 0.5 + 0.2 \cdot 1 = 0.4$, agent~$6$ follows
          agent~$1$, agents $4, 5$ at $0.5 \cdot 0.9 + 0.5 \cdot 0.4 = 0.65$.
    \item (d): $x_2(0) = 1$, consensus $0.4 \cdot 1 + 0.2 \cdot 1 = 0.6$, agents $4, 5$ at
          $0.5 \cdot 0.9 + 0.5 \cdot 0.6 = 0.75$. In the augmented model the closed group
          adds a seventh eigenvalue $1$.
  \end{itemize}
  \detailbuttons{\go{bk-aux-table}{Summary}\ \ret{m-aux}}
\end{frame}

\begin{frame}[label=bk-aux-table]{Auxiliary network: Table 2}
  \centering
  \scriptsize
  \renewcommand{\arraystretch}{1.25}
  \begin{tabular}{@{}c c >{\raggedright\arraybackslash}p{2.3cm} c >{\raggedright\arraybackslash}p{4.2cm} >{\raggedright\arraybackslash}p{3.6cm}@{}}
    \toprule
    Panel & $W$ & Prejudiced agent & P-dependent & Outcome & What it shows \\
    \midrule
    (a) & $W_{\mathrm{p}}$ & $2$, inside the closed group & all ($r = n$)
        & asymptotically stable (Corollary~3.3),
          $x(\infty) = u_2 \mathbf{1}_6 = 0.5\, \mathbf{1}_6$
        & the periodic cycle plays no role \\
    \midrule
    (b) & $W_{\mathrm{p}}$ & $5$, outside the closed group & $4, 5$ ($r = 2$)
        & $W^{22}$ not regular: not convergent (Theorem~3.2)
        & aperiodicity is needed again, for the P-independent agents \\
    \midrule
    (c) & $W_{\mathrm{a}}$ & $5$, outside the closed group & $4, 5$ ($r = 2$)
        & $W^{22}$ regular: convergent (Theorem~3.2), not asymptotically stable; agents
          $1,2,3,6$ at $0.4$, agents $4,5$ at $0.65$
        & convergence without asymptotic stability (Corollary~3.4) \\
    \midrule
    (d) & $W_{\mathrm{a}}$ & as (c), $x_2(0) = 1$, $u$ unchanged & $4, 5$ ($r = 2$)
        & as (c), but agents $1,2,3,6$ at $0.6$, agents $4,5$ at $0.75$
        & the limit depends on the initial opinions of the P-independent agents \\
    \bottomrule
  \end{tabular}
  \detailbuttons{\ret[Back to blocks and $V$]{bk-aux}\ \ret{m-aux}}
\end{frame}

\begin{frame}[label=bk-lemma-numbers]{The numbers of Figure 13}
  \small
  \begin{itemize}
    \item $W = A(1/2)$: strongly connected with a self-loop, primitive (Corollary~1.12),
          hence fully regular, as Lemma~3.5 requires. $\Lambda = \alpha I$,
          $\alpha \in [0, 0.995]$ (the lemma takes $\alpha < 1$), $u = x(0)$.
    \item $\alpha = 0$: $V = I$, so $c_0 = \frac16 \mathbf{1}_6$ (dotted line $1/6$).
    \item $\alpha \to 1$: $c_\alpha \to w$, the social power of the FDG model with
          $A = A(1/2)$: $w_1 = 2/7$, $w_2 = \dots = w_6 = 1/7$ (dashed lines), from
          $w(t)$ at $t = 1/2$.
    \item On the directed cycle $c_2$ first goes below $w_2$: its computed minimum is about
          $0.137$, at $\alpha \approx 0.65$.
    \item $c_\alpha$ exists on every network: $\rho(\alpha W) = \alpha < 1$ for every
          stochastic $W$, also where $w$ does not exist.
  \end{itemize}
  \detailbuttons{\ret{m-lemma}}
\end{frame}

% ---------------------------------------------------------------------------
% Conclusions and references
% ---------------------------------------------------------------------------

\begin{frame}[label=bk-table]{Comparison (Table 3)}
  \centering
  \scriptsize
  \renewcommand{\arraystretch}{0.95}
  \begin{tabular}{@{}>{\raggedright\arraybackslash}p{2.4cm}
                  >{\raggedright\arraybackslash}p{5.7cm}
                  >{\raggedright\arraybackslash}p{6.0cm}@{}}
    \toprule
    & FDG model (1) & FJ model (4) \\
    \midrule
    Resistance to the others
      & self-weight $a_{ii}$: inertia, a weight on the current opinion
      & attachment $1 - \lambda_i$ to the prejudice $u_i$: an anchor to a fixed value \\
    Stubborn agent (Def.~2.5)
      & $a_{ii} = 1$: a condition on the matrix only
      & $\lambda_i = 0$ and $u_i = x_i(0)$: a condition on the matrix and on the input \\
    Opinion of a stubborn agent & a state & an input \\
    Convergence
      & Table~1: sinks of $C(G)$ and their aperiodicity
      & all agents P-dependent (Corollary~3.3), or $W^{22}$ regular (Theorem~3.2) \\
    Asymptotic stability
      & never: the eigenvalue $1$ lies on the unit circle (Theorem~1.6)
      & if and only if all agents are P-dependent, i.e.\ $\rho(\Lambda W) < 1$
        (Corollary~3.3) \\
    Final opinions
      & consensus $(w^\top x(0))\mathbf{1}_n$ (Theorem~2.2), or disagreement with at least
        two aperiodic sinks (Theorem~2.4)
      & if all agents are P-dependent, $x(\infty) = Vu$ (9): convex combinations of the
        prejudices, independent of $x(0)$ \\
    Strongly connected network
      & consensus or no convergence, never disagreement
      & disagreement whenever the prejudices that count differ \\
    Aperiodicity & always needed & needed only for the P-independent agents \\
    Measure of power
      & social power $w$ (2): requires a consensus
      & influence centrality $c$ (12): requires $\rho(\Lambda W) < 1$ and $u = x(0)$ \\
    From one model to the other
      & with stubborn agents, an FJ model with $\lambda_i = 0$, $u_i = x_i(0)$ for the
        stubborn agents [5], or, removing them, an FJ model in which their opinions are
        prejudices (Figure~8)
      & an FDG model on the augmented graph (11), with $n$ virtual stubborn agents \\
    \bottomrule
  \end{tabular}
  \detailbuttons{\ret{m-table}}
\end{frame}

\begin{frame}[label=bk-refs]{References}
  \small
  \begin{enumerate}
    \item[{[1]}] Giacomo Baggio. DT averaging systems. Lecture slides, Network Dynamical
          Systems, University of Padova, 2025.
    \item[{[2]}] Giacomo Baggio. Elements of Graph Theory. Lecture slides, Network
          Dynamical Systems, University of Padova, 2025.
    \item[{[3]}] Giacomo Baggio. Elements of Matrix Theory. Lecture slides, Network
          Dynamical Systems, University of Padova, 2025.
    \item[{[4]}] Francesco Bullo. Lectures on Network Systems. Kindle Direct Publishing,
          1.7 edition, 2024. With contributions by J.\ Cort\'es, F.\ D\"orfler, and
          S.\ Mart\'inez.
    \item[{[5]}] Anton V.\ Proskurnikov and Roberto Tempo. A tutorial on modeling and
          analysis of dynamic social networks. Part I. Annual Reviews in Control,
          43:65--79, 2017.
  \end{enumerate}
\end{frame}
